\begin{definition}[Group]
    Let $G$ be a set. The structure $(G, \cdot)$ is a \emph{group} if it satisfies the following axioms:
    \begin{itemize}
        \item Closure: For all $a, b \in G$, $a \cdot b \in G$.
        \item Associativity: For all $a, b, c \in G$, $(a \cdot b) \cdot c = a \cdot (b \cdot c)$.
        \item Identity: There exists an element $e \in G$ such that for all $a \in G$, $e \cdot a = a \cdot e = a$.
        \item Inverse: For all $a \in G$, there exists an element $a^{-1} \in G$ such that $a \cdot a^{-1} = a^{-1} \cdot a = e$.
    \end{itemize}
    If the group also satisfies the commutative property (i.e. $a \cdot b = b \cdot a$ for all $a, b \in G$), then it is called an \emph{abelian group}.
\end{definition}
Commonly, we call $(G, +)$ an \emph{additive} group and $(G, \cdot)$ a \emph{multiplicative} group. In this case, the identity element are denoted by $0$ and $1$, respectively.

\begin{definition}[Field]
    Let $K$ be a set. The structure $(K, +, \cdot)$ is a \emph{field} if it satisfies the following properties:
    \begin{itemize}
        \item Additive group: $(K, +)$ is an abelian group.
        \item Multiplicative group: $(K\setminus\{0\}, \cdot)$ is an abelian group.
        \item Distributivity: For all $a, b, c \in K$, $a \cdot (b + c) = a \cdot b + a \cdot c$ and $(a + b) \cdot c = a \cdot c + b \cdot c$.
    \end{itemize}
\end{definition}

\begin{definition}[Characteristic]
    Let $K$ be a field. The \emph{characteristic} of $K$ is the smallest positive integer $n$ such that
    \begin{equation*}
        \underbrace{1 + \cdots + 1}_\text{$n$ times} = 0
    \end{equation*}
\end{definition}

\begin{definition}[Order]
    Let $K$ be a field. The \emph{order} of $K$ is the amount of elements in $K$; if the order of $K$ is finite, it is also called a Galois field.
    A Galois field of order $q$, denoted as $\mathds{F}_q$ or $GF(q)$, exists if and only if $q=p^n$ for some prime $p$; all Galois fields of order $q$ are isomorphic.
\end{definition}
Galois fields can be constructed in two ways:
\begin{itemize}
    \item Primitive fields: $GF(p) \cong \mathds{Z}_p$, with $p$ prime, is a field of order and characteristic $p$.
    \item $GF(p^n) = GF(p)[X]/P(X)$ with $P(X)$ irreducible polynomial in $GF(p)[X]$ of degree $n$,
        is a field of order $p^n$ and characteristic $p$.
\end{itemize}

\begin{definition}[Field extension]
    Let $K$ and $L$ be fields. If $K \subseteq L$, then we say that $K$ is a \emph{subfield} of $L$; conversely, 
    we say that $L$ is a \emph{field extension} of $K$.
    Given a field extension, denoted as $L/K$, the field $L$ is a vector space over $K$; the dimension of the vector space
    defines the \emph{degree of the extension}, and is denoted by $[L:K]$. In the case of finite fields, 
    $[GF(p^n):GF(p)] = n$.
\end{definition}

\begin{example}[Binary fields]
    The smallest Galois field is $GF(2)$, which is isomorphic to $\mathds{Z}_2 = \{0,1\}$.
    We can build $\mathds{F}_8 = GF(2^3)$ as $GF(2)[X]/\langle X^3+X+1\rangle$.
    The elements of $\mathds{F}_8$ are then the binary polynomials of degree strictly less than $n=3$; to avoid confusion, 
    the elements are written as polynomials in the variable $\alpha$, where $P(\alpha) = 0$.\\
    We have that $\mathds{F}_8 = \{0, 1, \alpha, 1+\alpha, \alpha^2, 1+\alpha^2, \alpha + \alpha^2, 1+\alpha+\alpha^2\}$, so the order of the field is $8$.
    Moreover, since $1+1=0$ in $GF(2)$, we have that $1+1=0$ in $\mathds{F}_8$ as well, so the characteristic of $\mathds{F}_8$ is $2$.\\
    By construction, $\mathds{F}_8$ is a field extension of $GF(2)$ of degree $3$; the vector space $\mathds{F}_8$ is thus a vector space over $GF(2)$ of dimension $3$,
    having as basis the elements $\{1, \alpha, \alpha^2\}$.
\end{example}
